\documentclass[aspectratio=169]{beamer}

% ─── Tema y colores ───────────────────────────────────────────
\usetheme{Madrid}
\usecolortheme{default}

\definecolor{azuloscuro}{RGB}{15, 23, 42}
\definecolor{azulmedio}{RGB}{30, 58, 138}
\definecolor{azulclaro}{RGB}{96, 165, 250}
\definecolor{verde}{RGB}{16, 185, 129}
\definecolor{naranja}{RGB}{245, 158, 11}
\definecolor{gris}{RGB}{100, 116, 139}
\definecolor{rojo}{RGB}{239, 68, 68}
\definecolor{fondoclaro}{RGB}{241, 245, 249}

\setbeamercolor{palette primary}{bg=azuloscuro, fg=white}
\setbeamercolor{palette secondary}{bg=azulmedio, fg=white}
\setbeamercolor{palette tertiary}{bg=azulclaro, fg=azuloscuro}
\setbeamercolor{structure}{fg=azulmedio}
\setbeamercolor{frametitle}{bg=azuloscuro, fg=white}
\setbeamercolor{title}{fg=white}
\setbeamercolor{subtitle}{fg=azulclaro}
\setbeamercolor{block title}{bg=azulmedio, fg=white}
\setbeamercolor{block body}{bg=fondoclaro, fg=azuloscuro}
\setbeamercolor{block title alerted}{bg=rojo, fg=white}
\setbeamercolor{block body alerted}{bg=red!10, fg=azuloscuro}
\setbeamercolor{block title example}{bg=verde, fg=white}
\setbeamercolor{block body example}{bg=green!10, fg=azuloscuro}

% ─── Paquetes ─────────────────────────────────────────────────
\usepackage[utf8]{inputenc}
\usepackage[spanish]{babel}
\usepackage{listings}
\usepackage{xcolor}
\usepackage{tikz}
\usetikzlibrary{positioning}
\usepackage{graphicx}
\usepackage{booktabs}
\usepackage{array}

% ─── Estilo de codigo ─────────────────────────────────────────
\lstdefinestyle{bash}{
  language=bash,
  basicstyle=\ttfamily\footnotesize\color{white},
  keywordstyle=\color{azulclaro}\bfseries,
  commentstyle=\color{gris}\itshape,
  stringstyle=\color{verde},
  backgroundcolor=\color{azuloscuro},
  frame=single,
  framerule=0pt,
  numbers=none,
  breaklines=true,
  showstringspaces=false,
  xleftmargin=5pt,
  xrightmargin=5pt,
}

\lstdefinestyle{yaml}{
  basicstyle=\ttfamily\tiny\color{white},
  keywordstyle=\color{azulclaro}\bfseries,
  commentstyle=\color{gris}\itshape,
  stringstyle=\color{verde},
  backgroundcolor=\color{azuloscuro},
  frame=single,
  framerule=0pt,
  numbers=none,
  breaklines=true,
  showstringspaces=false,
  xleftmargin=5pt,
  xrightmargin=5pt,
}

% ─── Metadatos ────────────────────────────────────────────────
\title{Integracion y Despliegue Continuo}
\subtitle{Git, GitHub y GitHub Actions para proyectos de procesamiento satelital}
\author{ADAIP}
\institute{Area de Desarrollos Avanzados de Imagenes y Percepcion}
\date{\today}

\setbeamertemplate{navigation symbols}{}
\setbeamertemplate{footline}[frame number]

% ══════════════════════════════════════════════════════════════
\begin{document}

% ─── Portada ──────────────────────────────────────────────────
{
\setbeamercolor{background canvas}{bg=azuloscuro}
\begin{frame}[plain]
  \vspace{1cm}
  \begin{center}
    {\color{azulclaro}\large\textbf{PROCESADORES SATELITALES}}\\[0.3cm]
    \rule{8cm}{0.4pt}\\[0.5cm]
    {\color{white}\Huge\textbf{Integracion y Despliegue}}\\[0.2cm]
    {\color{white}\Huge\textbf{Continuo}}\\[0.5cm]
    {\color{azulclaro}\large Git $\cdot$ GitHub $\cdot$ GitHub Actions}\\[1cm]
    \rule{8cm}{0.4pt}\\[0.5cm]
    {\color{gris}\small ADAIP --- \today}
  \end{center}
\end{frame}
}

% ─── Indice ───────────────────────────────────────────────────
\begin{frame}{Contenidos}
  \tableofcontents
\end{frame}

% ══════════════════════════════════════════════════════════════
\section{Evolucion del desarrollo de software}

\begin{frame}{Los anos 60-70: el codigo como obra artesanal}
  \begin{columns}
    \column{0.55\textwidth}
    \begin{block}{Como se trabajaba}
      \begin{itemize}
        \item El software lo escribia \textbf{una sola persona}
        \item No habia metodologia formal
        \item El codigo era considerado documentacion interna
        \item Los programas se entregaban en tarjetas perforadas
        \item No existia el concepto de ``version''
      \end{itemize}
    \end{block}

    \vspace{0.3cm}

    \begin{alertblock}{El problema}
      Si el programador se iba, el conocimiento se iba con el. No habia forma de saber que habia cambiado ni por que.
    \end{alertblock}

    \column{0.45\textwidth}
    \begin{center}
      \begin{tikzpicture}
        \node[fill=azuloscuro, text=white, rounded corners,
              minimum width=3.5cm, minimum height=1cm,
              text width=3cm, align=center] (p)
              {\small Programador\\unico};
        \node[fill=azulmedio!40, rounded corners,
              minimum width=3.5cm, minimum height=1cm,
              text width=3cm, align=center,
              below=0.6cm of p] (c)
              {\small Codigo en\\tarjetas perforadas};
        \node[fill=verde!30, rounded corners,
              minimum width=3.5cm, minimum height=1cm,
              text width=3cm, align=center,
              below=0.6cm of c] (m)
              {\small Maquina\\mainframe};
        \draw[->, thick] (p) -- (c);
        \draw[->, thick] (c) -- (m);
      \end{tikzpicture}
    \end{center}
  \end{columns}
\end{frame}

\begin{frame}{Los anos 80: equipos y el caos de la colaboracion}
  \begin{columns}
    \column{0.5\textwidth}
    \begin{block}{Que cambio}
      \begin{itemize}
        \item Los proyectos crecieron en complejidad
        \item Aparecieron los \textbf{equipos de desarrollo}
        \item Cada uno trabajaba en su propia copia del codigo
        \item La integracion se hacia al final del proyecto
        \item Aparecio el termino \textbf{``integration hell''}
      \end{itemize}
    \end{block}

    \column{0.5\textwidth}
    \begin{alertblock}{Integration Hell}
      Cuando varios desarrolladores trabajaban meses en paralelo y al intentar unir el codigo, nada funcionaba. Los cambios de uno rompian el trabajo del otro. Resolver conflictos llevaba semanas.
    \end{alertblock}
  \end{columns}

  \vspace{0.4cm}

  \begin{exampleblock}{Primera solucion: RCS y CVS (1982-1990)}
    Aparecieron los primeros sistemas de control de versiones. Un servidor central guardaba el historial. Solo una persona podia editar un archivo a la vez (\textbf{lock/unlock}).
  \end{exampleblock}
\end{frame}

\begin{frame}{Los anos 90: metodologias en cascada}
  \begin{center}
    \begin{tikzpicture}[node distance=0.3cm]
      \node[fill=azulmedio, text=white, rounded corners,
            minimum width=10cm, minimum height=0.7cm,
            text width=9cm, align=center] (req)
            {\small \textbf{1. Requisitos} (meses)};
      \node[fill=azulmedio!80, text=white, rounded corners,
            minimum width=8.5cm, minimum height=0.7cm,
            text width=8cm, align=center, below=of req] (dis)
            {\small \textbf{2. Diseno} (meses)};
      \node[fill=azulmedio!60, text=white, rounded corners,
            minimum width=7cm, minimum height=0.7cm,
            text width=6.5cm, align=center, below=of dis] (imp)
            {\small \textbf{3. Implementacion} (meses)};
      \node[fill=azulmedio!40, text=azuloscuro, rounded corners,
            minimum width=5.5cm, minimum height=0.7cm,
            text width=5cm, align=center, below=of imp] (test)
            {\small \textbf{4. Testing} (semanas)};
      \node[fill=azulmedio!20, text=azuloscuro, rounded corners,
            minimum width=4cm, minimum height=0.7cm,
            text width=3.5cm, align=center, below=of test] (dep)
            {\small \textbf{5. Deploy} (dias)};
    \end{tikzpicture}
  \end{center}

  \vspace{0.2cm}
  \begin{columns}
    \column{0.5\textwidth}
    \begin{alertblock}{El problema}
      El cliente recibia el producto \textbf{anos despues} de haberlo pedido. Para entonces, los requisitos habian cambiado y el software ya era obsoleto.
    \end{alertblock}

    \column{0.5\textwidth}
    \begin{block}{Avance: Subversion (SVN) en 2000}
      Permitia que multiples personas editaran el mismo archivo. El servidor central seguia siendo el unico punto de verdad.
    \end{block}
  \end{columns}
\end{frame}

\begin{frame}{Los 2000: el Manifiesto Agil y el cambio de paradigma}
  \begin{block}{Manifiesto Agil --- 2001}
    17 desarrolladores se reunieron en Utah y firmaron un documento que cambiaria la industria:
    \begin{itemize}
      \item \textbf{Individuos e interacciones} sobre procesos y herramientas
      \item \textbf{Software funcionando} sobre documentacion extensiva
      \item \textbf{Colaboracion con el cliente} sobre negociacion contractual
      \item \textbf{Respuesta al cambio} sobre seguir un plan
    \end{itemize}
  \end{block}

  \vspace{0.3cm}

  \begin{columns}
    \column{0.5\textwidth}
    \begin{exampleblock}{Consecuencias practicas}
      \begin{itemize}
        \item Ciclos cortos de entrega (\textbf{sprints} de 2 semanas)
        \item Feedback continuo del cliente
        \item Equipos pequenos y autonomos
        \item Testing como parte del desarrollo
      \end{itemize}
    \end{exampleblock}

    \column{0.5\textwidth}
    \begin{exampleblock}{Git nace en 2005}
      Linus Torvalds crea Git para manejar el desarrollo del kernel de Linux. A diferencia de SVN, cada desarrollador tiene una copia completa del historial. No hay servidor central obligatorio.
    \end{exampleblock}
  \end{columns}
\end{frame}

\begin{frame}{Los 2010: DevOps y la fusion de desarrollo y operaciones}
  \begin{columns}
    \column{0.5\textwidth}
    \begin{block}{El problema anterior}
      \begin{itemize}
        \item \textbf{Developers}: quieren cambiar cosas rapido
        \item \textbf{Operations}: quieren estabilidad, no cambios
        \item Resultado: conflicto permanente y deploys dolorosos
        \item \textit{``Funciona en mi maquina''} era el lema
      \end{itemize}
    \end{block}

    \vspace{0.3cm}

    \begin{exampleblock}{La solucion: DevOps}
      Unificar desarrollo y operaciones en un mismo equipo con las mismas herramientas y responsabilidades.
    \end{exampleblock}

    \column{0.5\textwidth}
    \begin{center}
      \begin{tikzpicture}
        \node[fill=azulmedio!30, rounded corners,
              minimum width=2.8cm, minimum height=1cm,
              text width=2.5cm, align=center] (dev)
              {\small \textbf{Dev}\\Desarrollo};
        \node[fill=verde!30, rounded corners,
              minimum width=2.8cm, minimum height=1cm,
              text width=2.5cm, align=center,
              right=0.5cm of dev] (ops)
              {\small \textbf{Ops}\\Operaciones};
        \node[fill=naranja!40, rounded corners,
              minimum width=3cm, minimum height=1cm,
              text width=2.8cm, align=center,
              below=0.8cm of dev, xshift=1.65cm] (devops)
              {\small \textbf{DevOps}\\CI/CD + Docker\\+ Automatizacion};
        \draw[->, thick] (dev) -- (devops);
        \draw[->, thick] (ops) -- (devops);
      \end{tikzpicture}
    \end{center}

    \begin{block}{Herramientas clave}
      Docker (2013), GitHub Actions (2019), Kubernetes (2014)
    \end{block}
  \end{columns}
\end{frame}

\begin{frame}{La linea del tiempo completa}
  \begin{center}
    \begin{tikzpicture}[scale=0.85]
      % Linea de tiempo
      \draw[thick, azulmedio] (0,0) -- (12,0);

      % Puntos y etiquetas
      \foreach \x/\y/\label/\sub in {
        0/0.3/1960/Programador unico,
        2/0.3/1982/RCS - primeras versiones,
        4/0.3/1994/SVN - servidor central,
        6/0.3/2001/Manifiesto Agil,
        8/0.3/2005/Git distribuido,
        10/0.3/2013/Docker,
        12/0.3/2019/GitHub Actions
      }{
        \node[fill=azulmedio, circle, minimum size=6pt, inner sep=0pt] at (\x,0) {};
        \node[above, text width=1.6cm, align=center, font=\tiny\bfseries] at (\x,\y) {\label};
        \node[below, text width=1.6cm, align=center, font=\tiny] at (\x,-0.3) {\sub};
      }
    \end{tikzpicture}
  \end{center}

  \vspace{0.4cm}

  \begin{alertblock}{La conclusion}
    Cada etapa surgio como respuesta a los problemas de la anterior. CI/CD no es una moda --- es la solucion natural a decadas de problemas de colaboracion, integracion y despliegue.
  \end{alertblock}
\end{frame}

% ══════════════════════════════════════════════════════════════
\section{Que es la Integracion Continua?}

\begin{frame}{El problema que resuelve}
  \begin{columns}
    \column{0.5\textwidth}
    \begin{alertblock}{Sin CI/CD}
      \begin{itemize}
        \item Cada desarrollador trabaja aislado
        \item La integracion se hace al final
        \item Los errores aparecen tarde
        \item El deploy es un proceso manual y riesgoso
        \item \textit{``Funciona en mi maquina''}
      \end{itemize}
    \end{alertblock}

    \column{0.5\textwidth}
    \begin{exampleblock}{Con CI/CD}
      \begin{itemize}
        \item El codigo se integra continuamente
        \item Los errores aparecen inmediatamente
        \item Los tests se corren automaticamente
        \item El deploy es automatico y reproducible
        \item El mismo codigo en todos los entornos
      \end{itemize}
    \end{exampleblock}
  \end{columns}
\end{frame}

\begin{frame}{Definiciones}
  \begin{block}{CI --- Integracion Continua}
    Cada vez que un desarrollador sube codigo, se ejecuta automaticamente una bateria de tests para verificar que nada esta roto. Si algo falla, se notifica inmediatamente antes de que el problema se propague.
  \end{block}

  \vspace{0.4cm}

  \begin{block}{CD --- Despliegue Continuo}
    Si los tests pasan, el codigo se despliega automaticamente al servidor correspondiente. Sin intervencion manual, sin pasos olvidados, siempre reproducible.
  \end{block}

  \vspace{0.4cm}

  \begin{center}
    \begin{tikzpicture}[node distance=0.6cm]
      \node[fill=azulmedio, text=white, rounded corners, minimum width=1.8cm, minimum height=0.6cm] (c) {Commit};
      \node[fill=verde, text=white, rounded corners, minimum width=1.8cm, minimum height=0.6cm, right=of c] (t) {Tests};
      \node[fill=naranja, text=white, rounded corners, minimum width=1.8cm, minimum height=0.6cm, right=of t] (b) {Build};
      \node[fill=rojo!80, text=white, rounded corners, minimum width=1.8cm, minimum height=0.6cm, right=of b] (d) {Deploy};
      \draw[->, thick, azulclaro] (c) -- (t);
      \draw[->, thick, azulclaro] (t) -- (b);
      \draw[->, thick, azulclaro] (b) -- (d);
    \end{tikzpicture}
  \end{center}
\end{frame}

% -------------------------------------------------------
\begin{frame}{Docker --- El entorno reproducible}
  \begin{block}{El problema que resuelve}
    \textit{``Funciona en mi maquina''} --- el codigo corre perfecto
    en el entorno del desarrollador pero falla en produccion por
    diferencias en versiones de Python, librerias, o configuracion
    del sistema operativo.
  \end{block}

  \vspace{0.3cm}

  \begin{columns}
    \column{0.5\textwidth}
    \begin{exampleblock}{Que es Docker}
      \begin{itemize}
        \item Empaqueta la aplicacion junto con
              \textbf{todo su entorno}: Python, librerias,
              configuracion
        \item El paquete resultante se llama
              \textbf{imagen}
        \item Al ejecutar una imagen se crea un
              \textbf{contenedor}
        \item Corre igual en cualquier maquina
      \end{itemize}
    \end{exampleblock}

    \column{0.5\textwidth}
    \begin{center}
      \begin{tikzpicture}[node distance=0.5cm]
        \node[fill=azulmedio!30, rounded corners,
              minimum width=3.5cm, minimum height=0.7cm,
              text width=3cm, align=center] (app)
              {\small Codigo Python};
        \node[fill=verde!30, rounded corners,
              minimum width=3.5cm, minimum height=0.7cm,
              text width=3cm, align=center,
              below=of app] (dep)
              {\small Dependencias\\(numpy, rasterio...)};
        \node[fill=naranja!30, rounded corners,
              minimum width=3.5cm, minimum height=0.7cm,
              text width=3cm, align=center,
              below=of dep] (so)
              {\small Sistema operativo\\(Ubuntu slim)};
        \node[fill=azuloscuro, text=white, rounded corners,
              minimum width=3.5cm, minimum height=0.7cm,
              text width=3cm, align=center,
              below=of so] (img)
              {\small Imagen Docker};
        \draw[->, thick] (app) -- (dep);
        \draw[->, thick] (dep) -- (so);
        \draw[->, thick] (so) -- (img);
      \end{tikzpicture}
    \end{center}
  \end{columns}

  \vspace{0.3cm}
  \begin{alertblock}{En el pipeline CI/CD}
    GitHub Actions construye la imagen Docker, la sube a GHCR,
    y Render la descarga y corre. El mismo paquete en todos los entornos.
  \end{alertblock}
\end{frame}
% ══════════════════════════════════════════════════════════════
\section{Git --- Control de versiones}

\begin{frame}{Que es Git?}
  Git es un sistema de \textbf{control de versiones distribuido}. Registra todos los cambios del codigo a lo largo del tiempo, permitiendo:

  \vspace{0.4cm}
  \begin{columns}
    \column{0.33\textwidth}
    \begin{center}
      \textbf{Historial completo}\\[0.2cm]
      {\small Cada cambio queda registrado con autor, fecha y mensaje}
    \end{center}

    \column{0.33\textwidth}
    \begin{center}
      \textbf{Ramas paralelas}\\[0.2cm]
      {\small Multiples lineas de desarrollo simultaneas}
    \end{center}

    \column{0.33\textwidth}
    \begin{center}
      \textbf{Reversibilidad}\\[0.2cm]
      {\small Volver a cualquier estado anterior del codigo}
    \end{center}
  \end{columns}

  \vspace{0.4cm}
  \begin{block}{Conceptos clave}
    \begin{itemize}
      \item \textbf{Repositorio}: carpeta con todo el historial del proyecto
      \item \textbf{Commit}: foto del estado del codigo en un momento dado
      \item \textbf{Rama (branch)}: linea de desarrollo independiente
      \item \textbf{Remote}: copia del repositorio en un servidor (GitHub)
    \end{itemize}
  \end{block}
\end{frame}

\begin{frame}[fragile]{Comandos fundamentales --- Parte 1}
  \begin{columns}
    \column{0.5\textwidth}
      \begin{exampleblock}{\texttt{git init}}
        Crea un repositorio en la carpeta actual
        \begin{lstlisting}[style=bash]
git init              # repo nuevo
git init nombre-repo  # con nombre
        \end{lstlisting}
      \end{exampleblock}

      \vspace{0.3cm}

      \begin{exampleblock}{\texttt{git add}}
        Prepara archivos para el proximo commit
        \begin{lstlisting}[style=bash]
git add archivo.py   # un archivo
git add .            # todos los cambios
git add src/         # una carpeta
        \end{lstlisting}
      \end{exampleblock}

    \column{0.5\textwidth}
      \begin{exampleblock}{\texttt{git commit}}
        Guarda los cambios con un mensaje
        \begin{lstlisting}[style=bash]
git commit -m "descripcion"
git commit --amend   # modificar el ultimo
        \end{lstlisting}
      \end{exampleblock}

      \vspace{0.3cm}

      \begin{exampleblock}{\texttt{git push}}
        Sube commits al servidor remoto
        \begin{lstlisting}[style=bash]
git push origin main
git push origin develop
git push --force origin main
        \end{lstlisting}
      \end{exampleblock}
  \end{columns}
\end{frame}

% -------------------------------------------------------
\begin{frame}{Flujo basico --- de lo local al remoto}
  \begin{center}
    \begin{tikzpicture}[node distance=0.6cm]
      \node[fill=blue!20, rounded corners,
            minimum width=3.5cm, minimum height=0.8cm] (init)
            {Carpeta vacia};
      \node[fill=gris!30, rounded corners,
            minimum width=3.5cm, minimum height=0.8cm,
            below=of init] (wd) {Working Directory};
      \node[fill=naranja!30, rounded corners,
            minimum width=3.5cm, minimum height=0.8cm,
            below=of wd] (sa) {Staging Area};
      \node[fill=verde!30, rounded corners,
            minimum width=3.5cm, minimum height=0.8cm,
            below=of sa] (repo) {Repositorio local};
      \node[fill=azuloscuro!80, text=white, rounded corners,
            minimum width=3.5cm, minimum height=0.8cm,
            below=of repo] (remote) {Remote (GitHub)};

      \draw[->, thick] (init)   -- node[right]{\small \texttt{git init}}   (wd);
      \draw[->, thick] (wd)     -- node[right]{\small \texttt{git add}}    (sa);
      \draw[->, thick] (sa)     -- node[right]{\small \texttt{git commit}} (repo);
      \draw[->, thick] (repo)   -- node[right]{\small \texttt{git push}}   (remote);
    \end{tikzpicture}
  \end{center}
\end{frame}

% -------------------------------------------------------
\begin{frame}[fragile]{Comandos fundamentales --- Parte 2}
  \begin{columns}
    \column{0.5\textwidth}
      \begin{exampleblock}{\texttt{git pull}}
        Descarga y fusiona cambios remotos
        \begin{lstlisting}[style=bash]
git pull origin main
git pull --rebase origin main
        \end{lstlisting}
      \end{exampleblock}

      \vspace{0.3cm}

      \begin{exampleblock}{\texttt{git fetch}}
        Descarga cambios SIN fusionarlos
        \begin{lstlisting}[style=bash]
git fetch origin
git fetch --all
        \end{lstlisting}
      \end{exampleblock}

    \column{0.5\textwidth}
      \begin{block}{fetch vs pull}
        \begin{itemize}
          \item \texttt{pull} = \texttt{fetch} + \texttt{merge}
          \item \texttt{fetch} permite revisar antes de fusionar
          \item \texttt{fetch} es mas seguro en ramas compartidas
        \end{itemize}
      \end{block}
  \end{columns}
\end{frame}

% -------------------------------------------------------
\begin{frame}{Flujo inverso --- del remoto a lo local}
  \begin{center}
    \begin{tikzpicture}[node distance=0.6cm]
      \node[fill=azuloscuro!80, text=white, rounded corners,
            minimum width=3.5cm, minimum height=0.8cm] (remote)
            {Remote (GitHub)};
      \node[fill=verde!30, rounded corners,
            minimum width=3.5cm, minimum height=0.8cm,
            below=of remote] (repo) {Repositorio local};
      \node[fill=gris!30, rounded corners,
            minimum width=3.5cm, minimum height=0.8cm,
            below=of repo] (wd) {Working Directory};

      % fetch: remote -> repo local
      \draw[->, thick, naranja] (remote.west)
            to[out=180, in=180]
            node[left]{\small \texttt{git fetch}}
            (repo.west);

      % pull: remote -> working directory (directo)
      \draw[->, thick, verde!70!black] (remote.east)
            to[out=0, in=0]
            node[right]{\small \texttt{git pull}}
            (wd.east);

      % merge: repo -> wd
      \draw[->, thick, dashed] (repo)
            -- node[right]{\small \texttt{git merge}}
            (wd);
    \end{tikzpicture}
  \end{center}

  \vspace{0.3cm}
  \begin{block}{Diferencia clave}
    \texttt{fetch} trae los cambios al repositorio local sin tocar tu codigo.
    \texttt{pull} los trae y los fusiona directamente en tu rama activa.
  \end{block}
\end{frame}

\begin{frame}[fragile]{Comandos de ramas}
  \begin{columns}
    \column{0.5\textwidth}
      \begin{exampleblock}{\texttt{git branch}}
        Crea y lista ramas
        \begin{lstlisting}[style=bash]
git branch            # listar ramas
git branch nueva-rama # crear rama
git branch -d rama    # eliminar rama
        \end{lstlisting}
      \end{exampleblock}

      \vspace{0.2cm}

      \begin{exampleblock}{\texttt{git checkout / git switch}}
        Cambia entre ramas
        \begin{lstlisting}[style=bash]
git checkout main
git checkout -b feature  # crear y cambiar
git switch develop
git switch -c nueva-rama
        \end{lstlisting}
      \end{exampleblock}

      \vspace{0.2cm}

      \begin{exampleblock}{\texttt{git merge}}
        Fusiona una rama en otra
        \begin{lstlisting}[style=bash]
git checkout main
git merge develop
        \end{lstlisting}
      \end{exampleblock}

    \column{0.5\textwidth}
      \begin{exampleblock}{\texttt{git status / log}}
        \begin{lstlisting}[style=bash]
git status
git log --oneline
git log --oneline --graph
git diff
        \end{lstlisting}
      \end{exampleblock}

      \vspace{0.2cm}

      \begin{exampleblock}{\texttt{git stash}}
        Guarda cambios temporalmente
        \begin{lstlisting}[style=bash]
git stash
git stash pop
git stash list
        \end{lstlisting}
      \end{exampleblock}

      \vspace{0.2cm}

      \begin{block}{Buena practica}
        Nunca trabajar directamente en \texttt{main}.
        Siempre crear una rama \texttt{feature/xxx}
        y fusionar via Pull Request.
      \end{block}
  \end{columns}
\end{frame}

\begin{frame}[fragile]{Tags --- Versiones del proyecto}
  Los tags marcan puntos especificos en la historia, tipicamente versiones de lanzamiento.

  \vspace{0.3cm}

  \begin{block}{Versionado semantico: MAJOR.MINOR.PATCH}
    \begin{center}
      \texttt{v\textcolor{rojo}{1}.\textcolor{naranja}{2}.\textcolor{verde}{3}}\\[0.3cm]
      \textcolor{rojo}{\textbf{MAJOR}}: cambio que rompe compatibilidad\\
      \textcolor{naranja}{\textbf{MINOR}}: funcionalidad nueva, compatible\\
      \textcolor{verde}{\textbf{PATCH}}: correccion de bugs
    \end{center}
  \end{block}

  \vspace{0.3cm}

  \begin{columns}
    \column{0.5\textwidth}
    \begin{lstlisting}[style=bash]
# Crear tag
git tag v1.0.0
git tag -a v1.0.0 -m "Primera version"

# Subir a GitHub
git push origin v1.0.0
git push origin --tags
    \end{lstlisting}

    \column{0.5\textwidth}
    \begin{lstlisting}[style=bash]
# Ver tags
git tag
git tag -l "v1.*"

# Ir a un tag especifico
git checkout v1.0.0
    \end{lstlisting}
  \end{columns}
\end{frame}

% ══════════════════════════════════════════════════════════════
\section{GitHub --- Trabajo en equipo}

\begin{frame}{GitHub vs Git}
  \begin{columns}
    \column{0.5\textwidth}
    \begin{block}{Git}
      \begin{itemize}
        \item Herramienta de control de versiones
        \item Corre en tu maquina local
        \item No requiere internet
        \item Maneja el historial del codigo
      \end{itemize}
    \end{block}

    \column{0.5\textwidth}
    \begin{block}{GitHub}
      \begin{itemize}
        \item Plataforma web para alojar repositorios
        \item Permite colaboracion entre equipos
        \item Issues, Projects, Actions, Releases
        \item Alternativas: GitLab, Bitbucket
      \end{itemize}
    \end{block}
  \end{columns}

  \vspace{0.5cm}

  \begin{alertblock}{Analogia}
    Git es el auto. GitHub es la ruta por donde circula.
  \end{alertblock}
\end{frame}

\begin{frame}{Estrategia de ramas para equipos}
  \begin{center}
    \begin{tikzpicture}[scale=0.9]
      \draw[thick, azulmedio] (0,2) -- (10,2);
      \node[left] at (0,2) {\small\textbf{main}};
      \node[fill=azulmedio, circle, minimum size=6pt, inner sep=0pt] at (0,2) {};
      \node[fill=azulmedio, circle, minimum size=6pt, inner sep=0pt] at (5,2) {};
      \node[fill=azulmedio, circle, minimum size=6pt, inner sep=0pt] at (10,2) {};
      \node[above] at (5,2.1) {\tiny v1.0};
      \node[above] at (10,2.1) {\tiny v1.1};

      \draw[thick, verde] (0,1) -- (10,1);
      \node[left] at (0,1) {\small\textbf{develop}};

      \draw[thick, naranja] (1,1) -- (2,0.2) -- (4,0.2) -- (5,1);
      \node[below] at (2.5,0.2) {\tiny feature/proc-sar};

      \draw[thick, naranja!70] (3,1) -- (3.5,0.5) -- (6,0.5) -- (7,1);
      \node[below] at (5,0.5) {\tiny feature/calibracion};

      \draw[dashed, azulclaro, thick] (5,1) -- (5,2);
      \draw[dashed, azulclaro, thick] (10,1) -- (10,2);
    \end{tikzpicture}
  \end{center}

  \vspace{0.2cm}
  \begin{itemize}
    \item \textbf{main}: codigo en produccion, siempre estable
    \item \textbf{develop}: integracion de funcionalidades, base para testing
    \item \textbf{feature/xxx}: desarrollo de funcionalidades individuales
  \end{itemize}
\end{frame}

\begin{frame}{Issues --- Gestion de tareas}
  Los \textbf{Issues} son el sistema de GitHub para gestionar tareas, reportar errores, proponer mejoras y discutir cambios dentro de un repositorio.

  \vspace{0.3cm}

  \begin{columns}
    \column{0.5\textwidth}
    \begin{block}{Tipos de Issues}
      \begin{itemize}
        \item \textbf{Bug}: algo no funciona correctamente
        \item \textbf{Feature}: nueva funcionalidad
        \item \textbf{Research}: investigacion tecnica
        \item \textbf{Docs}: documentacion pendiente
      \end{itemize}
    \end{block}

    \column{0.5\textwidth}
    \begin{block}{Buenas practicas}
      \begin{itemize}
        \item Titulo claro y descriptivo
        \item Descripcion con contexto
        \item Asignar responsable
        \item Agregar etiquetas (labels)
        \item Vincular con un Milestone
      \end{itemize}
    \end{block}
  \end{columns}

  \vspace{0.3cm}

  \begin{exampleblock}{Ejemplo para procesamiento SAR}
    \textbf{\#42 --- Implementar correccion radiometrica para Sentinel-1}\\
    \small Aplicar LUT de calibracion segun documentacion ESA. Incluir test con escena de referencia de la Patagonia.
  \end{exampleblock}
\end{frame}

\begin{frame}{Milestones y Projects}
  \begin{columns}
    \column{0.5\textwidth}
    \begin{block}{Milestones --- Hitos}
      Agrupan Issues relacionados con una meta o fecha limite.

      \vspace{0.2cm}
      Ejemplo:
      \begin{itemize}
        \item \textbf{v1.0 --- Procesador basico}\\{\small Issues \#1, \#2, \#5, \#8}
        \item \textbf{v1.1 --- Correccion atmosferica}\\{\small Issues \#9, \#12, \#15}
        \item \textbf{v2.0 --- Interfaz web}\\{\small Issues \#20, \#21, \#25}
      \end{itemize}
    \end{block}

    \column{0.5\textwidth}
    \begin{block}{Projects --- Tablero Kanban}
      Vista visual del estado de todas las tareas.

      \vspace{0.2cm}
      \begin{tabular}{|c|c|c|}
        \hline
        \textbf{Backlog} & \textbf{En curso} & \textbf{Listo} \\
        \hline
        \#15 & \#8  & \#1 \\
        \#20 & \#12 & \#2 \\
        \#21 &      & \#5 \\
        \hline
      \end{tabular}
    \end{block}
  \end{columns}

  \vspace{0.3cm}
  \begin{alertblock}{Flujo recomendado}
    Issue $\rightarrow$ asignar a Milestone $\rightarrow$ mover en Project board $\rightarrow$ cerrar con PR
  \end{alertblock}
\end{frame}

\begin{frame}{Pull Requests --- Revision de codigo}
  Un \textbf{Pull Request (PR)} es una solicitud para fusionar una rama en otra.

  \vspace{0.3cm}

  \begin{enumerate}
    \item Desarrollador crea rama \texttt{feature/mi-funcionalidad}
    \item Trabaja y hace commits
    \item Abre un PR hacia \texttt{develop}
    \item Otro miembro del equipo revisa el codigo
    \item Se discuten cambios, se hacen ajustes
    \item Si los tests pasan y hay aprobacion $\rightarrow$ se fusiona
  \end{enumerate}

  \vspace{0.3cm}

  \begin{exampleblock}{Ventajas}
    \begin{itemize}
      \item Nadie sube codigo directamente a \texttt{main}
      \item Revision obligatoria antes de integrar
      \item Historial de decisiones tecnicas documentado
      \item Los tests corren automaticamente antes de aprobar
    \end{itemize}
  \end{exampleblock}
\end{frame}

% ══════════════════════════════════════════════════════════════
\section{GitHub Actions --- Automatizacion}

\begin{frame}{Que son los Workflows?}
  Un \textbf{workflow} es un proceso automatizado definido en un archivo YAML dentro del repositorio.

  \vspace{0.3cm}

  \begin{block}{Caracteristicas}
    \begin{itemize}
      \item Viven en \texttt{.github/workflows/nombre.yml}
      \item Se activan por eventos: push, pull request, horario, manual
      \item Corren en maquinas virtuales de GitHub (runners)
      \item Son codigo: se versionan junto al proyecto
    \end{itemize}
  \end{block}

  \vspace{0.3cm}

  \begin{columns}
    \column{0.33\textwidth}
    \begin{center}
      \textbf{Disparador}\\[0.2cm]
      {\small push, pull request, schedule}
    \end{center}

    \column{0.33\textwidth}
    \begin{center}
      \textbf{Jobs}\\[0.2cm]
      {\small test, build, deploy}
    \end{center}

    \column{0.33\textwidth}
    \begin{center}
      \textbf{Steps}\\[0.2cm]
      {\small comandos individuales}
    \end{center}
  \end{columns}
\end{frame}

\begin{frame}[fragile]{Anatomia de un workflow}
  \begin{lstlisting}[style=yaml]
name: CI/CD                      # nombre del workflow

on:                              # cuando se ejecuta
  push:
    branches: [main, develop]    # en push a estas ramas
  pull_request:
    branches: [main]             # en PR hacia main

jobs:                            # trabajos a ejecutar

  test:                          # ID del job
    name: Correr tests           # nombre visible en GitHub
    runs-on: ubuntu-latest       # sistema operativo del runner
    steps:
      - name: Descargar codigo
        uses: actions/checkout@v4

      - name: Instalar Python
        uses: actions/setup-python@v5
        with:
          python-version: "3.13"

      - name: Correr tests
        run: pytest tests/ -v
  \end{lstlisting}
\end{frame}

\begin{frame}[fragile]{Pipeline completo: test $\to$ build $\to$ deploy}
  \begin{lstlisting}[style=yaml]
jobs:

  test:                            # Job 1: tests
    runs-on: ubuntu-latest
    steps:
      - uses: actions/checkout@v4
      - run: pip install -r requirements.txt
      - run: pytest tests/ -v

  build:                           # Job 2: solo si test paso
    needs: test
    runs-on: ubuntu-latest
    steps:
      - uses: actions/checkout@v4
      - uses: docker/build-push-action@v5
        with:
          push: true
          tags: ghcr.io/mi-usuario/mi-app:latest

  deploy:                          # Job 3: solo en main
    needs: [test, build]
    if: github.ref == 'refs/heads/main'
    steps:
      - run: curl -X POST "${{ secrets.DEPLOY_HOOK }}"
  \end{lstlisting}
\end{frame}

\begin{frame}{Conceptos clave de los workflows}
  \begin{columns}
    \column{0.5\textwidth}
    \begin{block}{Disparadores (on:)}
      \begin{itemize}
        \item \texttt{push} --- al subir codigo
        \item \texttt{pull\_request} --- al abrir un PR
        \item \texttt{schedule} --- con cron (periodico)
        \item \texttt{workflow\_dispatch} --- manual
        \item \texttt{release} --- al publicar una version
      \end{itemize}
    \end{block}

    \vspace{0.3cm}

    \begin{block}{Secrets}
      Variables sensibles guardadas en GitHub Settings. Se referencian como \texttt{\$\{\{ secrets.NOMBRE \}\}} y nunca aparecen en los logs.
    \end{block}

    \column{0.5\textwidth}
    \begin{block}{Contextos disponibles}
      \begin{itemize}
        \item \texttt{github.ref} --- rama actual
        \item \texttt{github.actor} --- quien hizo push
        \item \texttt{github.repository} --- usuario/repo
        \item \texttt{github.sha} --- hash del commit
      \end{itemize}
    \end{block}

    \vspace{0.3cm}

    \begin{block}{needs --- Dependencias}
      \texttt{needs: test} hace que un job espere a que otro termine exitosamente antes de ejecutarse.
    \end{block}
  \end{columns}
\end{frame}

% ══════════════════════════════════════════════════════════════
\section{Ejemplo practico --- Procesador SAR}

\begin{frame}{Caso de uso: procesador SAR con CI/CD}
  \begin{block}{Escenario}
    Equipo de 3 personas desarrollando un procesador de imagenes Sentinel-1 en Python. Necesitan garantizar que cada cambio no rompa el procesamiento existente.
  \end{block}

  \vspace{0.3cm}

  \begin{columns}
    \column{0.5\textwidth}
    \begin{exampleblock}{Estructura del proyecto}
      \begin{itemize}
        \item \texttt{processor/} --- algoritmos SAR
        \item \texttt{tests/} --- tests con imagenes de referencia
        \item \texttt{data/} --- escenas de calibracion
        \item \texttt{Dockerfile} --- entorno reproducible
        \item \texttt{.github/workflows/}
      \end{itemize}
    \end{exampleblock}

    \column{0.5\textwidth}
    \begin{exampleblock}{Flujo de trabajo}
      \begin{enumerate}
        \item Issue \#42: nueva funcion de filtrado
        \item Rama \texttt{feature/filtro-lee}
        \item Desarrollo + tests unitarios
        \item PR hacia \texttt{develop}
        \item CI corre tests automaticamente
        \item Revision del equipo
        \item Merge $\rightarrow$ deploy a staging
        \item Validacion $\rightarrow$ merge a \texttt{main}
      \end{enumerate}
    \end{exampleblock}
  \end{columns}
\end{frame}

\begin{frame}[fragile]{Workflow para procesamiento satelital}
  \begin{lstlisting}[style=yaml]
name: CI --- Procesador SAR

on:
  push:
    branches: [main, develop]
  pull_request:
    branches: [develop]

jobs:

  test:
    name: Tests de procesamiento
    runs-on: ubuntu-latest
    steps:
      - uses: actions/checkout@v4

      - name: Instalar dependencias
        run: |
          pip install numpy scipy rasterio pytest
          pip install -r requirements.txt

      - name: Tests unitarios
        run: pytest tests/unit/ -v

      - name: Tests con escenas de referencia
        run: pytest tests/integration/ -v --timeout=120

      - name: Calidad del codigo
        run: |
          flake8 processor/
          black processor/ --check
  \end{lstlisting}
\end{frame}

\begin{frame}{Beneficios concretos para el equipo}
  \begin{columns}
    \column{0.5\textwidth}
    \begin{alertblock}{Sin CI/CD}
      \begin{itemize}
        \item El procesador SAR funciona solo en la PC del Dr. Garcia
        \item Al pasar a produccion falla por version diferente de GDAL
        \item Nadie sabe que cambio rompio el coregistro
        \item El deploy tarda 2 dias y requiere intervencion manual
      \end{itemize}
    \end{alertblock}

    \column{0.5\textwidth}
    \begin{exampleblock}{Con CI/CD}
      \begin{itemize}
        \item Docker garantiza el mismo entorno siempre
        \item Si un cambio rompe el procesamiento, se detecta en minutos
        \item El historial de Git muestra exactamente que cambio
        \item El deploy es automatico y tarda 5 minutos
      \end{itemize}
    \end{exampleblock}
  \end{columns}

  \vspace{0.5cm}

  \begin{alertblock}{Principio fundamental}
    \centering
    El codigo que paso los tests en CI es exactamente el mismo que se despliega en produccion.
  \end{alertblock}
\end{frame}

% ══════════════════════════════════════════════════════════════
\section{Caso de estudio completo --- FastApi2}

\begin{frame}{Arquitectura del sistema FastApi2}
  Sistema web completo con CI/CD, compuesto por dos repositorios independientes.

  \vspace{0.3cm}

  \begin{center}
    \begin{tikzpicture}[node distance=0.8cm]
      \node[fill=azulclaro!30, rounded corners, minimum width=3cm, minimum height=0.8cm] (user)
            {\small Usuario (navegador)};
      \node[fill=verde!30, rounded corners, minimum width=3cm, minimum height=0.8cm,
            right=1.cm of user] (front)
            {\small React (Frontend)};
      \node[fill=azulmedio!40, rounded corners, minimum width=3cm, minimum height=0.8cm,
            right=1.cm of front] (back)
            {\small FastAPI (Backend)};
      \node[fill=naranja!30, rounded corners, minimum width=3cm, minimum height=0.8cm,
            right=1.cm of back] (db)
            {\small PostgreSQL};

      \draw[->, thick] (user) -- node[above]{\tiny HTTP} (front);
      \draw[->, thick] (front) -- node[above]{\tiny JSON} (back);
      \draw[->, thick] (back) -- node[above]{\tiny SQL} (db);
    \end{tikzpicture}
  \end{center}

  \vspace{0.3cm}

  \begin{columns}
    \column{0.5\textwidth}
    \begin{block}{Repositorio backend}
      \texttt{sebaheredia/FastApi2}\\[0.2cm]
      \begin{itemize}
        \item \texttt{main.py} --- endpoints REST
        \item \texttt{database.py} --- conexion a PostgreSQL
        \item \texttt{models.py} --- tabla users (SQLAlchemy)
        \item \texttt{schemas.py} --- validacion JSON (Pydantic)
        \item \texttt{conftest.py} --- setup de tests
        \item \texttt{test\_main.py} --- tests con pytest
        \item \texttt{Dockerfile} --- imagen del backend
      \end{itemize}
    \end{block}

    \column{0.5\textwidth}
    \begin{block}{Repositorio frontend}
      \texttt{sebaheredia/fastapi2-frontend}\\[0.2cm]
      \begin{itemize}
        \item \texttt{src/App.js} --- componente React principal
        \item \texttt{src/api.js} --- llamadas al backend
        \item \texttt{src/App.css} --- estilos
        \item \texttt{nginx.conf} --- servidor web
        \item \texttt{Dockerfile} --- multi-stage Node/nginx
      \end{itemize}
    \end{block}
  \end{columns}
\end{frame}

\begin{frame}[fragile]{Docker --- Backend}
  El backend usa un Dockerfile de una sola etapa. Crea las tablas en PostgreSQL antes de arrancar el servidor:

  \vspace{0.3cm}

  \begin{lstlisting}[style=yaml]
FROM python:3.13-slim
WORKDIR /app

# Copiar deps primero para aprovechar cache de Docker
COPY requirements.txt .
RUN pip install --no-cache-dir -r requirements.txt

COPY . .
EXPOSE 8000

# Dos comandos en secuencia con &&:
# 1. Crear tablas en PostgreSQL (create_all)
# 2. Solo si tuvo exito, arrancar uvicorn
CMD ["sh", "-c",
  "python -c 'from database import engine; import models;
   models.Base.metadata.create_all(bind=engine)'
   && uvicorn main:app --host 0.0.0.0 --port $PORT"]
  \end{lstlisting}

  \begin{alertblock}{Por que psycopg2-binary es critico}
    Sin este driver en \texttt{requirements.txt}, la imagen se buildea sin soporte para PostgreSQL y la app cae silenciosamente a SQLite aunque \texttt{DATABASE\_URL} este correctamente configurada.
  \end{alertblock}
\end{frame}

\begin{frame}[fragile]{Docker --- Frontend (multi-stage)}
  El frontend usa un build de dos etapas: Node compila React y nginx sirve los archivos estaticos:

  \begin{lstlisting}[style=yaml]
# Etapa 1: Node compila React
FROM node:20-slim AS builder
WORKDIR /app
COPY package.json ./
RUN npm install
COPY . .

# La URL del backend se embebe en el bundle en tiempo de build
ARG REACT_APP_API_URL
ENV REACT_APP_API_URL=$REACT_APP_API_URL
RUN npm run build     # genera archivos estaticos en /build

# Etapa 2: nginx sirve los archivos estaticos
FROM nginx:alpine
COPY --from=builder /app/build /usr/share/nginx/html
COPY nginx.conf /etc/nginx/conf.d/default.conf
EXPOSE 80
CMD ["nginx", "-g", "daemon off;"]
  \end{lstlisting}

  \vspace{0.2cm}
  \begin{exampleblock}{Ventaja del multi-stage}
    La imagen final pesa $\sim$25 MB (nginx) en lugar de $\sim$300 MB (Node). Node solo se usa para compilar, no para servir.
  \end{exampleblock}
\end{frame}

\begin{frame}{Base de datos PostgreSQL en Render}
  \begin{columns}
    \column{0.5\textwidth}
    \begin{alertblock}{SQLite dentro del Docker (incorrecto)}
      \begin{itemize}
        \item Datos guardados en un archivo dentro del contenedor
        \item Al redesplegar, Render destruye el contenedor
        \item El archivo desaparece con el contenedor
        \item Todos los datos se pierden en cada deploy
      \end{itemize}
    \end{alertblock}

    \column{0.5\textwidth}
    \begin{exampleblock}{PostgreSQL externo (correcto)}
      \begin{itemize}
        \item Servidor independiente en Render
        \item El contenedor puede destruirse y recrearse
        \item La DB nunca se toca en los deploys
        \item Los datos persisten siempre
      \end{itemize}
    \end{exampleblock}
  \end{columns}

  \vspace{0.4cm}

  \begin{block}{Conexion via variable de entorno}
    Render inyecta \texttt{DATABASE\_URL} dentro del contenedor al arrancarlo. El codigo la lee con \texttt{os.getenv()}:

    \vspace{0.2cm}
    \texttt{DATABASE\_URL = os.getenv("DATABASE\_URL", "sqlite:///./database.db")}

    \vspace{0.1cm}
    {\small Si la variable existe $\rightarrow$ usa PostgreSQL. Si no $\rightarrow$ usa SQLite (desarrollo local).}
  \end{block}
\end{frame}

\begin{frame}{GHCR --- Registro de imagenes Docker}
  GitHub Container Registry almacena las imagenes construidas por GitHub Actions. Render las descarga directamente.

  \vspace{0.3cm}

  \begin{center}
    \begin{tikzpicture}[node distance=0.7cm]
      \node[fill=azulmedio!30, rounded corners, minimum width=2.8cm, minimum height=0.7cm] (ga)
            {\small GitHub Actions};
      \node[fill=verde!30, rounded corners, minimum width=2.8cm, minimum height=0.7cm,
            right=1cm of ga] (ghcr)
            {\small GHCR};
      \node[fill=naranja!30, rounded corners, minimum width=2.8cm, minimum height=0.7cm,
            right=1cm of ghcr] (render)
            {\small Render};

      \draw[->, thick] (ga) -- node[above]{\tiny push} (ghcr);
      \draw[->, thick] (ghcr) -- node[above]{\tiny pull} (render);
    \end{tikzpicture}
  \end{center}

  \vspace{0.3cm}

  \begin{columns}
    \column{0.5\textwidth}
    \begin{block}{Imagenes del backend}
      \begin{itemize}
        \item \texttt{ghcr.io/sebaheredia/fastapi2:main}
        \item \texttt{ghcr.io/sebaheredia/fastapi2:develop}
      \end{itemize}
    \end{block}

    \column{0.5\textwidth}
    \begin{block}{Imagenes del frontend}
      \begin{itemize}
        \item \texttt{ghcr.io/sebaheredia/fastapi2-frontend:main}
        \item \texttt{ghcr.io/sebaheredia/fastapi2-frontend:develop}
      \end{itemize}
    \end{block}
  \end{columns}

  \vspace{0.3cm}
  \begin{exampleblock}{Ventaja clave}
    La imagen que corre en produccion es exactamente la misma que fue testeada en CI. No hay diferencias de entorno.
  \end{exampleblock}
\end{frame}

\begin{frame}[fragile]{Pipeline CI/CD del backend}
  \begin{lstlisting}[style=yaml]
name: CI/CD Backend
on:
  push:
    branches: [main, develop]

jobs:
  test:                              # Job 1: tests
    runs-on: ubuntu-latest
    steps:
      - uses: actions/checkout@v4
      - uses: actions/setup-python@v5
        with: {python-version: "3.13"}
      - run: pip install -r requirements.txt
      - run: pytest test_main.py -v  # si falla, se detiene

  docker:                            # Job 2: build imagen
    needs: test
    steps:
      - uses: docker/login-action@v3
        with: {registry: ghcr.io, ...}
      - uses: docker/build-push-action@v5
        with:
          push: true
          tags: ghcr.io/sebaheredia/fastapi2:main

  deploy:                            # Job 3: notificar a Render
    needs: docker
    if: github.ref == 'refs/heads/main'
    run: |
      curl -X PATCH "api.render.com/v1/services/.../image" ...
      curl -X POST "api.render.com/v1/services/.../deploys" ...
  \end{lstlisting}
\end{frame}

\begin{frame}[fragile]{Pipeline CI/CD del frontend}
  El frontend tiene una particularidad: \texttt{REACT\_APP\_API\_URL} debe determinarse antes del build.

  \begin{lstlisting}[style=yaml]
jobs:
  docker:
    steps:
      - uses: actions/checkout@v4
      - uses: docker/login-action@v3
        with: {registry: ghcr.io, ...}

      # Determinar URL del backend segun la rama
      - name: Definir URL del backend
        id: backend_url
        run: |
          if [ "${{ github.ref }}" = "refs/heads/main" ]; then
            echo "url=${{ secrets.BACKEND_PRODUCTION_URL }}" >> $GITHUB_OUTPUT
          else
            echo "url=${{ secrets.BACKEND_STAGING_URL }}" >> $GITHUB_OUTPUT
          fi

      # La URL se embebe en el bundle de React durante el build
      - uses: docker/build-push-action@v5
        with:
          build-args: |
            REACT_APP_API_URL=${{ steps.backend_url.outputs.url }}
          push: true
          tags: ghcr.io/sebaheredia/fastapi2-frontend:main
  \end{lstlisting}
\end{frame}

\begin{frame}{Servicios desplegados en Render}
  \begin{block}{4 servicios web + 1 base de datos}
    \begin{center}
      \begin{tabular}{llll}
        \toprule
        \textbf{Servicio} & \textbf{Tipo} & \textbf{Rama} & \textbf{URL} \\
        \midrule
        fastapi2-staging & Image/GHCR & develop & ...staging-docker.onrender.com \\
        fastapi2-production & Image/GHCR & main & ...production-docker.onrender.com \\
        fastapi2-frontend-staging & Image/GHCR & develop & ...frontend-staging.onrender.com \\
        fastapi2-frontend-main & Image/GHCR & main & ...frontend-main.onrender.com \\
        fastapi2-db & PostgreSQL & --- & red interna de Render \\
        \bottomrule
      \end{tabular}
    \end{center}
  \end{block}

  \vspace{0.3cm}

  \begin{columns}
    \column{0.5\textwidth}
    \begin{block}{Secrets en GitHub (backend)}
      \begin{itemize}
        \item \texttt{RENDER\_API\_KEY}
        \item \texttt{RENDER\_STAGING\_SERVICE\_ID}
        \item \texttt{RENDER\_PRODUCTION\_SERVICE\_ID}
        \item \texttt{GHCR\_TOKEN}
      \end{itemize}
    \end{block}

    \column{0.5\textwidth}
    \begin{block}{Secrets en GitHub (frontend)}
      \begin{itemize}
        \item \texttt{RENDER\_API\_KEY}
        \item \texttt{RENDER\_FRONTEND\_STAGING\_SERVICE\_ID}
        \item \texttt{RENDER\_FRONTEND\_PRODUCTION\_SERVICE\_ID}
        \item \texttt{BACKEND\_STAGING\_URL}
        \item \texttt{BACKEND\_PRODUCTION\_URL}
      \end{itemize}
    \end{block}
  \end{columns}
\end{frame}

\begin{frame}{Flujo completo de un cambio en produccion}
  \begin{center}
    \begin{tikzpicture}[node distance=0.4cm, scale=0.85, every node/.style={scale=0.85}]
      \node[fill=azuloscuro, text=white, rounded corners, minimum width=3cm] (dev)
            {Desarrollador};
      \node[fill=azulmedio!50, rounded corners, minimum width=3cm, below=of dev] (push)
            {git push origin develop};
      \node[fill=verde!40, rounded corners, minimum width=3cm, below=of push] (ci)
            {GitHub Actions: tests};
      \node[fill=verde!40, rounded corners, minimum width=3cm, below=of ci] (build)
            {Build imagen Docker};
      \node[fill=naranja!40, rounded corners, minimum width=3cm, below=of build] (ghcr)
            {Push a GHCR};
      \node[fill=azulclaro!40, rounded corners, minimum width=3cm, below=of ghcr] (render)
            {Render descarga imagen};
      \node[fill=verde!60, rounded corners, minimum width=3cm, below=of render] (live)
            {Servicio staging activo};

      \draw[->, thick] (dev) -- (push);
      \draw[->, thick] (push) -- (ci);
      \draw[->, thick] (ci) -- node[right]{\tiny OK} (build);
      \draw[->, thick] (build) -- (ghcr);
      \draw[->, thick] (ghcr) -- (render);
      \draw[->, thick] (render) -- (live);

      \node[fill=rojo!40, rounded corners, right=1.5cm of ci, minimum width=2cm] (fail)
            {Falla: stop};
      \draw[->, thick, rojo] (ci) -- node[above]{\tiny Error} (fail);
    \end{tikzpicture}
  \end{center}
\end{frame}

\begin{frame}[fragile]{CORS --- Comunicacion entre dominios}
  \begin{block}{El problema}
    Frontend (\texttt{fastapi2-frontend.onrender.com}) y backend (\texttt{fastapi2-backend.onrender.com}) estan en dominios distintos. El navegador bloquea los requests entre dominios por seguridad.
  \end{block}

  \vspace{0.3cm}

  \begin{exampleblock}{La solucion: middleware CORS en FastAPI}
    \ttfamily\small
    from fastapi.middleware.cors import CORSMiddleware\\[0.2cm]
    app = FastAPI()\\[0.2cm]
    app.add\_middleware(\\
    \quad CORSMiddleware,\\
    \quad allow\_origins=\textbf{[}"*"\textbf{]},\\
    \quad allow\_methods=\textbf{[}"*"\textbf{]},\\
    \quad allow\_headers=\textbf{[}"*"\textbf{]},\\
    )
    \normalfont
  \end{exampleblock}

  \begin{alertblock}{Importante}
    \texttt{add\_middleware} es una llamada de funcion, NO un decorador. No lleva \texttt{@} al principio.
  \end{alertblock}
\end{frame}

% ══════════════════════════════════════════════════════════════
\section{Resumen}

\begin{frame}{Resumen}
  \begin{columns}
    \column{0.5\textwidth}
    \begin{block}{Lo que vimos hoy}
      \begin{enumerate}
        \item \textbf{Git}: control de versiones
          \begin{itemize}
            \item add, commit, push, pull, fetch
            \item merge, checkout, stash
            \item tags y versionado semantico
          \end{itemize}
        \item \textbf{GitHub}: plataforma colaborativa
          \begin{itemize}
            \item Issues, Milestones, Projects
            \item Ramas: main, develop, feature
            \item Pull Requests
          \end{itemize}
        \item \textbf{GitHub Actions}
          \begin{itemize}
            \item Workflows, jobs, steps
            \item Disparadores y condiciones
            \item Secrets y contextos
          \end{itemize}
      \end{enumerate}
    \end{block}

    \column{0.5\textwidth}
    \begin{exampleblock}{Proximos pasos}
      \begin{enumerate}
        \item Crear repositorio para el procesador SAR
        \item Configurar ramas main/develop
        \item Escribir primeros tests automatizados
        \item Crear el primer workflow de CI
        \item Agregar Docker para reproducibilidad
        \item Configurar deploy automatico
      \end{enumerate}
    \end{exampleblock}

    \vspace{0.3cm}
    \begin{alertblock}{Herramientas necesarias}
      Git, VSCode, Docker Desktop, cuenta en GitHub
    \end{alertblock}
  \end{columns}
\end{frame}

% ─── Cierre ───────────────────────────────────────────────────
{
\setbeamercolor{background canvas}{bg=azuloscuro}
\begin{frame}[plain]
  \begin{center}
    \vspace{2cm}
    {\color{white}\Huge\textbf{Gracias!}}\\[0.5cm]
    \rule{6cm}{0.4pt}\\[0.5cm]
    {\color{gris}\large ADAIP}\\[0.2cm]
    {\color{gris}\small Area de Desarrollos Avanzados de Imagenes y Percepcion}\\[1cm]
    \rule{6cm}{0.4pt}\\[0.5cm]
    {\color{azulclaro}\small Repositorios del ejemplo practico:}\\[0.2cm]
    {\color{white}\small github.com/sebaheredia/FastApi2}\\
    {\color{white}\small github.com/sebaheredia/fastapi2-frontend}
  \end{center}
\end{frame}
}

\end{document}
